% GENERATED FILE. Edit canonical lessons, then run npm run generate:books.
% canonical-source-hash: fnv1a64:1bcae1f9672e009e
% canonical-lessons: RU-C178-trenirovatsya, RU-C178-begat, RU-C178-prygat, RU-C178-katatsya, RU-C178-zagorat

\chapter{To train, To run about, To jump, To go for a ride, To sunbathe}
\label{ch:ru-to-train-to-run-about-to-jump-to-go-for-a-ride-to-sunbathe}
\hlchaptermodality{178}

\begin{chapteropening}
\textbf{By the end of this chapter:} \emph{I can say to train, to run about, to jump, to go for a ride and to sunbathe.}

Complete the last lesson of chapter 178: I can say to train, to run about, to jump, to go for a ride and to sunbathe.
\end{chapteropening}

\section[\texorpdfstring{\ru{тренироваться} (trenirovátsya)}{тренироваться (trenirovátsya)}]{\ru{тренироваться} (trenirovátsya) — to train, to practise}

\label{lesson:RU-C178-trenirovatsya}



\begin{quote}
Before the new one: say the Russian for to quarrel, then the Russian for to make up.
\end{quote}

\subsection*{You'll want to know: \ru{тренироваться}}
\textbf{\ru{тренироваться}} (\emph{trenirovátsya}) — \textquotedblleft{}to train, to practise\textquotedblright{}.

\begin{grammarlens}[title={What you've built}]
One more word for this chapter.
\end{grammarlens}

\subsection*{Your turn}
\begin{itemize}
\raggedright
  \item \emph{Say it:} \emph{trenirovátsya}
  \item \emph{Say it:} \emph{trenirovátsya}, once more
  \item \emph{Say it:} say \emph{trenirovátsya}
  \item \emph{Recall:} say the Russian for to quarrel, then the Russian for to make up, then say \emph{trenirovátsya} again
\end{itemize}

\begin{tcolorbox}[breakable,colback=teal!4,colframe=teal!35!black,title={Before you move on}]
What does \ru{тренироваться} mean? (\textquotedblleft{}To train, to practise\textquotedblright{}.) Say it once more.
\end{tcolorbox}
\section[\texorpdfstring{\ru{бегать} (bégat)}{бегать (bégat)}]{\ru{бегать} (bégat) — to run about, to jog}

\label{lesson:RU-C178-begat}



\begin{quote}
Before the new one: say the Russian for to make up, then the Russian for to train.
\end{quote}

\subsection*{You'll want to know: \ru{бегать}}
\textbf{\ru{бегать}} (\emph{bégat}) — \textquotedblleft{}to run about, to jog\textquotedblright{}.

\begin{grammarlens}[title={What you've built}]
One more word for this chapter.
\end{grammarlens}

\subsection*{Your turn}
\begin{itemize}
\raggedright
  \item \emph{Say it:} \emph{bégat}
  \item \emph{Say it:} \emph{bégat}, once more
  \item \emph{Say it:} say \emph{bégat}
  \item \emph{Recall:} say the Russian for to make up, then the Russian for to train, then say \emph{bégat} again
\end{itemize}

\begin{tcolorbox}[breakable,colback=teal!4,colframe=teal!35!black,title={Before you move on}]
What does \ru{бегать} mean? (\textquotedblleft{}To run about, to jog\textquotedblright{}.) Say it once more.
\end{tcolorbox}
\section[\texorpdfstring{\ru{прыгать} (prýgat)}{прыгать (prýgat)}]{\ru{прыгать} (prýgat) — to jump}

\label{lesson:RU-C178-prygat}



\begin{quote}
Before the new one: say the Russian for to train, then the Russian for to run about.
\end{quote}

\subsection*{You'll want to know: \ru{прыгать}}
\textbf{\ru{прыгать}} (\emph{prýgat}) — \textquotedblleft{}to jump\textquotedblright{}.

\begin{grammarlens}[title={What you've built}]
One more word for this chapter.
\end{grammarlens}

\subsection*{Your turn}
\begin{itemize}
\raggedright
  \item \emph{Say it:} \emph{prýgat}
  \item \emph{Say it:} \emph{prýgat}, once more
  \item \emph{Say it:} say \emph{prýgat}
  \item \emph{Recall:} say the Russian for to train, then the Russian for to run about, then say \emph{prýgat} again
\end{itemize}

\begin{tcolorbox}[breakable,colback=teal!4,colframe=teal!35!black,title={Before you move on}]
What does \ru{прыгать} mean? (\textquotedblleft{}To jump\textquotedblright{}.) Say it once more.
\end{tcolorbox}
\section[\texorpdfstring{\ru{кататься} (katátsya)}{кататься (katátsya)}]{\ru{кататься} (katátsya) — to go for a ride, to skate}

\label{lesson:RU-C178-katatsya}



\begin{quote}
Before the new one: say the Russian for to run about, then the Russian for to jump.
\end{quote}

\subsection*{You'll want to know: \ru{кататься}}
\textbf{\ru{кататься}} (\emph{katátsya}) — \textquotedblleft{}to go for a ride, to skate\textquotedblright{}.

\begin{grammarlens}[title={What you've built}]
One more word for this chapter.
\end{grammarlens}

\subsection*{Your turn}
\begin{itemize}
\raggedright
  \item \emph{Say it:} \emph{katátsya}
  \item \emph{Say it:} \emph{katátsya}, once more
  \item \emph{Say it:} say \emph{katátsya}
  \item \emph{Recall:} say the Russian for to run about, then the Russian for to jump, then say \emph{katátsya} again
\end{itemize}

\begin{tcolorbox}[breakable,colback=teal!4,colframe=teal!35!black,title={Before you move on}]
What does \ru{кататься} mean? (\textquotedblleft{}To go for a ride, to skate\textquotedblright{}.) Say it once more.
\end{tcolorbox}
\section[\texorpdfstring{\ru{загорать} (zagorát)}{загорать (zagorát)}]{\ru{загорать} (zagorát) — to sunbathe}

\label{lesson:RU-C178-zagorat}



\begin{quote}
Before the new one: say the Russian for to jump, then the Russian for to go for a ride.
\end{quote}

\subsection*{You'll want to know: \ru{загорать}}
\textbf{\ru{загорать}} (\emph{zagorát}) — \textquotedblleft{}to sunbathe\textquotedblright{}.

\begin{grammarlens}[title={What you've built}]
One more word for this chapter.
\end{grammarlens}

\subsection*{Your turn}
\begin{itemize}
\raggedright
  \item \emph{Say it:} \emph{zagorát}
  \item \emph{Say it:} \emph{zagorát}, once more
  \item \emph{Say it:} say \emph{zagorát}
  \item \emph{Recall:} say the Russian for to jump, then the Russian for to go for a ride, then say \emph{zagorát} again
\end{itemize}

\begin{tcolorbox}[breakable,colback=teal!4,colframe=teal!35!black,title={Before you move on}]
What does \ru{загорать} mean? (\textquotedblleft{}To sunbathe\textquotedblright{}.) Say it once more.
\end{tcolorbox}
